%% Correction des données %%

Ces visualisations peuvent également contribuer à mettre en avant les irrégularités des bases de données. Dans le cas de l'INA, nous avons soulevé que certaines politiques documentaires ne circulaient qu'à l'oral ou par mails, ce qui signifie que nous ne pouvons pas les suivre. Grâce à la visualisation de ces collections, nous pouvons cibler ces irrégularités, comprendre pourquoi elles existent et mettre en place une médiation pour permettre aux utilisateurs d'accéder à ces données ou même les corriger. Il semble moins insurmontable de cibler des erreurs qui sont ostensiblement visibles plutôt que d'aller rechercher la trace de chaque petit changement de \gls{pd} dans des milliers de mails. Lors de la réalisation de tests de visualisation du \gls{td} chronologique de la chaîne \enquote{France 2}, nous avons notamment relevé des erreurs se prêtant à ce cas. La chronologie commence en 1992 parce c'est la date de création de la chaîne sous le nom \enquote{France 2}. Ce qui nous interpelle, c'est qu'avant 1995, date à laquelle le dépôt légal a été mis en place à l'INA, il y a un nombre plus important de notices, plus que toutes les autres années et toutes les notices sont en IDENT (Identification), donc sans aucun \gls{td}. Cela force la réflexion sur ce sujet. 
%REDEMANDER EXPLICATION. 

%% Attention aux biais %%

%% Les biais sont très courants dans les datavisualisations. On peut faire dire n'importe quoi à des images, en l'occurence des graphiques. C'est pourquoi il est important que les documentalistes restent impliqués dès le départ. 

%% Conclusion de la sous-partie 2 : %%

Cela dit il est important de souligner que cet outil a été pensé pour les bases de données du dépôt légal. Il n'est pas possible de faire des recherches dans les champs utilisés par le dépôt légal. De la même façon qu'il tente de régler des problèmes qui ont sans doute été réglés par les nouveaux outils à venir. Certains problèmes soulignés ont déjà été étudiés et corrigés par le service qualité de l'INA. Mais ces corrections concernent principalement la qualité des données. L'avantage de cet outil et plus précisément des datavisualisations est qu'il a été adapté pour correspondre à différentes bases de données. Ainsi, il est possible de l'utiliser sur d'autres bases, comme le lac de données car les enjeux d'accessibilité aux collections dues à cet enchaînement de politiques documentaires est propre à ces collections. Cet outil reste pertinent pour d'autres bases de données plus adaptées aux nouveaux modes de consommation des données, sans perdre de sa pertinence, même si il nécessite d'être adapté en conséquence. 

%% Conclusion du chapitre :%%

Cela dit, nous ne devons pas négliger que cet outil est à destination de personnes ayant pour objectif d'utiliser les bases de données de l'INA et plus précisément du dépôt légal. L'utilisateur peut rencontrer de nouveaux contenus par hasard, mais il reste important d'avoir une idée de départ de ce qu'il recherche : une chaîne, une genre etc... En partant d'une des clefs d'entrée de notre outil peut découvrir des documents qu'il n'aurait jamais rencontré involontairement lors d'une recherche sur Hyperbase mais il ne peut pas partir de vraiment rien. Cet outil a avant tout été imaginé en appui à la recherche, de ce fait, c'est la fonctionnalité à laquelle nous avons prêter le plus d'attention. 